\documentclass[11pt,a4paper]{article}

\usepackage[utf8]{inputenc}
\usepackage[T1]{fontenc}
\usepackage[spanish]{babel}
\usepackage{amsmath,amssymb,mathtools}
\usepackage{geometry}
\usepackage{lmodern}
\usepackage{microtype}

\geometry{margin=2.4cm}
\setlength{\parindent}{0pt}
\setlength{\parskip}{0.45em}

\newcommand{\R}{\mathbb{R}}
\newcommand{\sen}{\operatorname{sen}}

\begin{document}

\begin{center}
{\LARGE\bfseries Análisis Matemático II}\\[0.4em]
{\Large Tema 5: Teorema de la función implícita y de la función inversa}\\[0.4em]
{\large Soluciones}
\end{center}

\section{Ejercicio 1}
\subsection{Solución}

Sea:

\[
F(x,y)=x^2+y^2-2.
\]

En el punto \((1,1)\):

\[
F(1,1)=0.
\]

Además:

\[
F_y(x,y)=2y
\]

y:

\[
F_y(1,1)=2\neq0.
\]

Por tanto, la ecuación define implícitamente a:

\[
y=f(x)
\]

en un entorno de \((1,1)\).

La primera derivada es:

\[
y'
=
-\frac{F_x}{F_y}
=
-\frac{x}{y}.
\]

Así:

\[
\boxed{
y'(1)=-1
}
\]

Derivando:

\[
2x+2yy'=0,
\]

se obtiene:

\[
2+2(y')^2+2yy''=0.
\]

Por tanto:

\[
y''
=
-\frac{1+(y')^2}{y}.
\]

En \((1,1)\):

\[
\boxed{
y''(1)=-2
}
\]

\section{Ejercicio 2}
\subsection{Solución}

Sea:

\[
F(x,y)=y^3+x^3-6xy.
\]

La pendiente de la recta tangente es:

\[
y'
=
-\frac{F_x}{F_y}.
\]

Se tiene:

\[
F_x=3x^2-6y,
\]

\[
F_y=3y^2-6x.
\]

En \((3,3)\):

\[
F_x(3,3)=9,
\qquad
F_y(3,3)=9.
\]

Por tanto:

\[
y'(3)=-1.
\]

La recta tangente es:

\[
y-3=-(x-3).
\]

Así:

\[
\boxed{
y=-x+6
}
\]

\section{Ejercicio 3}
\subsection{Solución}

Sea:

\[
F(x,y)=xy+x+\ln y.
\]

En \((0,1)\):

\[
F(0,1)=0.
\]

Además:

\[
F_y=x+\frac1y.
\]

Por tanto:

\[
F_y(0,1)=1\neq0.
\]

Así, la ecuación define implícitamente a:

\[
y=f(x)
\]

en un entorno de \((0,1)\).

La derivada es:

\[
y'
=
-\frac{F_x}{F_y}
=
-\frac{y+1}{x+\frac1y}.
\]

En \((0,1)\):

\[
y'(0)=-2.
\]

La recta tangente es:

\[
y-1=-2x.
\]

Por tanto:

\[
\boxed{
y=1-2x
}
\]

\section{Ejercicio 4}
\subsection{Solución}

Sea:

\[
F(x,y)
=
(x^2+y^2)^2-4xy.
\]

Se tiene:

\[
F_x
=
4x(x^2+y^2)-4y,
\]

\[
F_y
=
4y(x^2+y^2)-4x.
\]

En \((1,1)\):

\[
F_x(1,1)=4,
\qquad
F_y(1,1)=4.
\]

Por tanto:

\[
y'(1)
=
-\frac{F_x(1,1)}{F_y(1,1)}
=
-1.
\]

La recta tangente es:

\[
y-1=-(x-1).
\]

Así:

\[
\boxed{
x+y=2
}
\]

\section{Ejercicio 5}
\subsection{Solución}

Sea:

\[
h(x,y)
=
x^2+y^3+xy+x^3+ay.
\]

En \((0,0)\):

\[
h(0,0)=0.
\]

Para definir:

\[
y=f(x),
\]

debe cumplirse:

\[
h_y(0,0)\neq0.
\]

Como:

\[
h_y
=
3y^2+x+a,
\]

se tiene:

\[
h_y(0,0)=a.
\]

Por tanto:

\[
\boxed{
a\neq0
}
\]

Para definir:

\[
x=g(y),
\]

debe cumplirse:

\[
h_x(0,0)\neq0.
\]

Pero:

\[
h_x
=
2x+y+3x^2
\]

y:

\[
h_x(0,0)=0
\]

para cualquier valor de \(a\).

Por tanto:

\[
\boxed{
\text{No define }x\text{ como función implícita diferenciable de }y
\text{ para ningún }a
}
\]

\section{Ejercicio 6}
\subsection{Solución}

Sea:

\[
F(x,y)
=
1-\ln\left(
1+\frac1{e^{xy}}
\right)
=
0.
\]

Entonces:

\[
\ln(1+e^{-xy})=1.
\]

Por tanto:

\[
e^{-xy}=e-1
\]

y:

\[
xy=-\ln(e-1).
\]

Como \(a\neq0\), localmente:

\[
y
=
-\frac{\ln(e-1)}{x}.
\]

Así:

\[
y'
=
\frac{\ln(e-1)}{x^2}.
\]

Usando:

\[
ab=-\ln(e-1),
\]

se obtiene:

\[
\boxed{
f'(a)
=
-\frac ba
}
\]

Derivando nuevamente:

\[
y''
=
-\frac{2\ln(e-1)}{x^3}.
\]

Por tanto:

\[
\boxed{
f''(a)
=
\frac{2b}{a^2}
}
\]

\section{Ejercicio 7}
\subsection{Solución}

Las circunferencias son:

\[
C_1:
(x-2)^2+y^2=2
\]

y:

\[
C_2:
x^2+(y-2)^2=2.
\]

Sus centros son:

\[
(2,0),
\qquad
(0,2),
\]

y ambas tienen radio:

\[
\sqrt2.
\]

La distancia entre los centros es:

\[
\sqrt{(2-0)^2+(0-2)^2}
=
2\sqrt2.
\]

Como coincide con la suma de los radios, son tangentes exteriormente en el punto medio de los centros:

\[
\boxed{
(1,1)
}
\]

Para:

\[
F_1(x,y)
=
(x-2)^2+y^2-2,
\]

\[
(F_1)_y(1,1)=2\neq0.
\]

Por tanto, \(C_1\) define localmente a \(y\) como función de \(x\), y:

\[
y_1'
=
-\frac{(F_1)_x}{(F_1)_y}.
\]

En \((1,1)\):

\[
y_1'(1)=1.
\]

Para:

\[
F_2(x,y)
=
x^2+(y-2)^2-2,
\]

\[
(F_2)_y(1,1)=-2\neq0.
\]

Así:

\[
y_2'
=
-\frac{(F_2)_x}{(F_2)_y}
\]

y:

\[
y_2'(1)=1.
\]

Las dos curvas tienen la misma recta tangente:

\[
y-1=x-1.
\]

Por tanto:

\[
\boxed{
y=x
}
\]

\section{Ejercicio 8}
\subsection{Solución}

Sea:

\[
F(x,y,z)
=
2x+3y-xe^{x+y-2z}-8.
\]

En:

\[
(-1,3,1),
\]

se tiene:

\[
F(-1,3,1)=0.
\]

Además:

\[
F_z
=
2xe^{x+y-2z}.
\]

Por tanto:

\[
F_z(-1,3,1)=-2\neq0.
\]

La ecuación define implícitamente:

\[
z=f(x,y).
\]

Las derivadas son:

\[
z_x
=
-\frac{F_x}{F_z},
\qquad
z_y
=
-\frac{F_y}{F_z}.
\]

Se tiene:

\[
F_x
=
2-(1+x)e^{x+y-2z},
\]

\[
F_y
=
3-xe^{x+y-2z}.
\]

En el punto:

\[
F_x(-1,3,1)=2,
\]

\[
F_y(-1,3,1)=4.
\]

Así:

\[
\boxed{
z_x(-1,3)=1
}
\]

\[
\boxed{
z_y(-1,3)=2
}
\]

\section{Ejercicio 9}
\subsection{Solución}

Sea:

\[
F(x,y,z)
=
yze^{x+3}
+
3xe^{z-y}.
\]

En:

\[
(1,-1,3),
\]

\[
F(1,-1,3)=0.
\]

Además:

\[
F_z
=
ye^{x+3}
+
3xe^{z-y}.
\]

Por tanto:

\[
F_z(1,-1,3)
=
2e^4\neq0.
\]

Así, la ecuación define:

\[
z=z(x,y)
\]

en un entorno de \((1,-1)\).

Calculamos:

\[
F_x
=
yze^{x+3}
+
3e^{z-y},
\]

\[
F_y
=
ze^{x+3}
-
3xe^{z-y}.
\]

En el punto:

\[
F_x(1,-1,3)=0,
\qquad
F_y(1,-1,3)=0.
\]

Por tanto:

\[
\boxed{
z_x(1,-1)=0
}
\]

\[
\boxed{
z_y(1,-1)=0
}
\]

\section{Ejercicio 10}
\subsection{Solución}

Sea:

\[
F(x,y,z)
=
1+x+y+z-e^{-(x+y+z)}.
\]

Se tiene:

\[
F_z
=
1+e^{-(x+y+z)}>0.
\]

Por tanto, en cualquier punto que satisfaga la ecuación:

\[
F(x,y,z)=0,
\]

puede definirse localmente:

\[
z=f(x,y).
\]

Además:

\[
F_x
=
F_y
=
F_z
=
1+e^{-(x+y+z)}.
\]

Por tanto:

\[
z_x
=
-\frac{F_x}{F_z}
=
-1,
\]

\[
z_y
=
-\frac{F_y}{F_z}
=
-1.
\]

Así:

\[
\boxed{
\frac{\partial z}{\partial x}=-1,
\qquad
\frac{\partial z}{\partial y}=-1
}
\]

y:

\[
\boxed{
dz=-dx-dy
}
\]

\section{Ejercicio 11}
\subsection{Solución}

Sea:

\[
z=x^2+y^2
\]

con la condición:

\[
F(x,y)
=
x^2+y^3+y
=
0.
\]

Como:

\[
F_y
=
3y^2+1>0,
\]

la ecuación define:

\[
y=y(x).
\]

Por tanto:

\[
y'
=
-\frac{F_x}{F_y}
=
-\frac{2x}{3y^2+1}.
\]

Derivando:

\[
z=x^2+y^2,
\]

se obtiene:

\[
\frac{\partial z}{\partial x}
=
2x+2yy'.
\]

Sustituyendo:

\[
\boxed{
\frac{\partial z}{\partial x}
=
2x
-
\frac{4xy}{3y^2+1}
}
\]

Equivalentemente:

\[
\boxed{
\frac{\partial z}{\partial x}
=
\frac{
2x(3y^2-2y+1)
}{
3y^2+1
}
}
\]

\section{Ejercicio 12}
\subsection{Solución}

Sea:

\[
F(x,y,z)
=
z^3+xz+y.
\]

En:

\[
(x,y,z)=(1,-2,1),
\]

se tiene:

\[
F(1,-2,1)=0.
\]

Además:

\[
F_z
=
3z^2+x.
\]

Por tanto:

\[
F_z(1,-2,1)=4\neq0.
\]

La ecuación define implícitamente:

\[
z=f(x,y).
\]

Las primeras derivadas son:

\[
z_x
=
-\frac{F_x}{F_z}
=
-\frac{z}{3z^2+x},
\]

\[
z_y
=
-\frac{F_y}{F_z}
=
-\frac1{3z^2+x}.
\]

Así:

\[
\boxed{
\nabla f(1,-2)
=
\left(
-\frac14,
-\frac14
\right)
}
\]

Para las segundas derivadas:

\[
z_{xx}
=
-
\frac{
F_{xx}
+
2F_{xz}z_x
+
F_{zz}z_x^2
}{
F_z
},
\]

\[
z_{xy}
=
-
\frac{
F_{xy}
+
F_{xz}z_y
+
F_{yz}z_x
+
F_{zz}z_xz_y
}{
F_z
},
\]

\[
z_{yy}
=
-
\frac{
F_{yy}
+
2F_{yz}z_y
+
F_{zz}z_y^2
}{
F_z
}.
\]

En el punto:

\[
F_{xz}=1,
\qquad
F_{zz}=6,
\]

y el resto de derivadas segundas de \(F\) son nulas.

Se obtiene:

\[
z_{xx}(1,-2)=\frac1{32},
\]

\[
z_{xy}(1,-2)=z_{yx}(1,-2)=-\frac1{32},
\]

\[
z_{yy}(1,-2)=-\frac3{32}.
\]

Por tanto:

\[
\boxed{
H_f(1,-2)
=
\frac1{32}
\begin{pmatrix}
1&-1\\
-1&-3
\end{pmatrix}
}
\]

\section{Ejercicio 13}
\subsection{Solución}

Sea:

\[
F(x,y,z)
=
z^2+z-xy-1.
\]

En \((1,1,1)\):

\[
F(1,1,1)=0.
\]

Además:

\[
F_z=2z+1
\]

y:

\[
F_z(1,1,1)=3\neq0.
\]

Por tanto, existe localmente:

\[
z=f(x,y).
\]

Las primeras derivadas son:

\[
z_x
=
\frac{y}{2z+1},
\]

\[
z_y
=
\frac{x}{2z+1}.
\]

En \((1,1)\):

\[
z_x(1,1)=z_y(1,1)=\frac13.
\]

Las segundas derivadas en el punto son:

\[
z_{xx}(1,1)
=
-\frac2{27},
\]

\[
z_{xy}(1,1)
=
z_{yx}(1,1)
=
\frac7{27},
\]

\[
z_{yy}(1,1)
=
-\frac2{27}.
\]

Tomando:

\[
h=x-1,
\qquad
k=y-1,
\]

el polinomio de Taylor de segundo orden es:

\[
\boxed{
P_2(x,y)
=
1
+
\frac{h+k}{3}
-
\frac{h^2+k^2}{27}
+
\frac{7hk}{27}
}
\]

es decir:

\[
\boxed{
\begin{aligned}
P_2(x,y)
={}&
1
+
\frac{x-1}{3}
+
\frac{y-1}{3}\\
&-
\frac{(x-1)^2}{27}
+
\frac{7(x-1)(y-1)}{27}
-
\frac{(y-1)^2}{27}.
\end{aligned}
}
\]

\section{Ejercicio 14}
\subsection{Solución}

Sea:

\[
F(x,y,z)
=
\cos\left(\frac{\pi z}{2}\right)
-xy+x+y-z^2.
\]

En \((1,1,1)\):

\[
F(1,1,1)=0.
\]

Además:

\[
F_z
=
-\frac{\pi}{2}
\sen\left(
\frac{\pi z}{2}
\right)
-2z.
\]

Así:

\[
F_z(1,1,1)
=
-\frac{\pi+4}{2}\neq0.
\]

Por tanto:

\[
z=z(x,y)
\]

está definida implícitamente en un entorno de \((1,1)\).

Como:

\[
F_x=-y+1,
\qquad
F_y=-x+1,
\]

se tiene:

\[
z_x(1,1)=z_y(1,1)=0.
\]

Además:

\[
F_{xy}=-1,
\]

y en el punto:

\[
z_{xx}=z_{yy}=0.
\]

Para la derivada mixta:

\[
z_{xy}
=
-\frac{F_{xy}}{F_z}
=
-\frac{2}{\pi+4}.
\]

Por tanto:

\[
\boxed{
P_2(x,y)
=
1
-
\frac{2}{\pi+4}
(x-1)(y-1)
}
\]

\section{Ejercicio 15}
\subsection{Solución}

Sea:

\[
f(x,y)
=
\left(
e^x+e^y,
e^x-e^y
\right).
\]

Su matriz jacobiana es:

\[
J_f(x,y)
=
\begin{pmatrix}
e^x&e^y\\
e^x&-e^y
\end{pmatrix}.
\]

Su determinante es:

\[
|J_f(x,y)|
=
-2e^{x+y}\neq0
\]

para todo \((x,y)\in\R^2\).

Por tanto:

\[
\boxed{
f\text{ es localmente invertible en todo punto de }\R^2
}
\]

La matriz jacobiana de la inversa es:

\[
J_{f^{-1}}
=
J_f^{-1}.
\]

Así:

\[
\boxed{
J_{f^{-1}}
=
\begin{pmatrix}
\dfrac1{2e^x}&\dfrac1{2e^x}\\[2mm]
\dfrac1{2e^y}&-\dfrac1{2e^y}
\end{pmatrix}
}
\]

Sea:

\[
u=e^x+e^y,
\qquad
v=e^x-e^y.
\]

Entonces:

\[
e^x=\frac{u+v}{2},
\qquad
e^y=\frac{u-v}{2}.
\]

Por tanto:

\[
\boxed{
f^{-1}(u,v)
=
\left(
\ln\frac{u+v}{2},
\ln\frac{u-v}{2}
\right)
}
\]

con:

\[
u+v>0,
\qquad
u-v>0.
\]

\section{Ejercicio 16}
\subsection{Solución}

Sea:

\[
f(x,y)
=
\left(
\frac{x}{y},
\frac{x^3+y^3}{x^3-y^3}
\right).
\]

Su dominio exige:

\[
y\neq0,
\qquad
x^3-y^3\neq0.
\]

La segunda componente puede escribirse en función de:

\[
u=\frac{x}{y}.
\]

En efecto:

\[
\frac{x^3+y^3}{x^3-y^3}
=
\frac{u^3+1}{u^3-1}.
\]

Por tanto, las dos componentes no son funcionalmente independientes.

El determinante jacobiano es:

\[
\boxed{
|J_f(x,y)|=0
}
\]

en todo su dominio.

Así:

\[
\boxed{
\text{No puede aplicarse el teorema de la función inversa en ningún punto}
}
\]

\section{Ejercicio 17}
\subsection{Solución}

Sea:

\[
f(x,y)
=
\left(
e^x\cos y,
e^x\sen y
\right).
\]

Su matriz jacobiana es:

\[
J_f(x,y)
=
\begin{pmatrix}
e^x\cos y&-e^x\sen y\\
e^x\sen y&e^x\cos y
\end{pmatrix}.
\]

Por tanto:

\[
|J_f(x,y)|
=
e^{2x}
\left(
\cos^2y+\sen^2y
\right)
=
e^{2x}>0.
\]

Así:

\[
\boxed{
f\text{ posee inversa local en todo punto de }\R^2
}
\]

Sin embargo:

\[
f(x,y+2\pi)=f(x,y).
\]

Por tanto:

\[
\boxed{
f\text{ no es inyectiva globalmente}
}
\]

Esto no contradice el teorema de la función inversa, ya que dicho teorema garantiza invertibilidad local, no global.

\section{Ejercicio 18}
\subsection{Solución}

Sea:

\[
f(x,y)
=
\left(
x+\alpha y,
e^x+\alpha e^y
\right).
\]

La matriz jacobiana es:

\[
J_f(x,y)
=
\begin{pmatrix}
1&\alpha\\
e^x&\alpha e^y
\end{pmatrix}.
\]

Su determinante es:

\[
|J_f(x,y)|
=
\alpha(e^y-e^x).
\]

Por tanto:

\[
|J_f(x,y)|\neq0
\]

si y sólo si:

\[
\alpha\neq0
\]

y:

\[
x\neq y.
\]

Así:

\[
\boxed{
f\text{ es localmente invertible en }(x,y)
\iff
\alpha\neq0
\text{ y }
x\neq y
}
\]

En particular:

\[
\boxed{
\alpha\neq0
}
\]

es necesario para que exista algún punto en el que pueda aplicarse el teorema de la función inversa.

\section{Ejercicio 19}
\subsection{Solución}

Sea:

\[
f(x,y,z)
=
\left(
e^x,
\sen(x+y),
e^z
\right).
\]

Su matriz jacobiana es:

\[
J_f(x,y,z)
=
\begin{pmatrix}
e^x&0&0\\
\cos(x+y)&\cos(x+y)&0\\
0&0&e^z
\end{pmatrix}.
\]

Por tanto:

\[
|J_f(x,y,z)|
=
e^{x+z}\cos(x+y).
\]

En \((0,0,0)\):

\[
|J_f(0,0,0)|=1\neq0.
\]

Así:

\[
\boxed{
f\text{ es localmente invertible en }(0,0,0)
}
\]

Las hipótesis del teorema no se cumplen cuando:

\[
\cos(x+y)=0.
\]

Es decir:

\[
\boxed{
x+y
=
\frac{\pi}{2}
+k\pi,
\qquad
k\in\mathbb Z
}
\]

En esos puntos:

\[
|J_f|=0.
\]

\section{Ejercicio 20}
\subsection{Solución}

Sea:

\[
f(x,y)
=
\left(
x^2-y^2,
2xy
\right).
\]

Su matriz jacobiana es:

\[
J_f(x,y)
=
\begin{pmatrix}
2x&-2y\\
2y&2x
\end{pmatrix}.
\]

El determinante es:

\[
|J_f(x,y)|
=
4(x^2+y^2).
\]

Por tanto, puede aplicarse el teorema de la función inversa cuando:

\[
\boxed{
(x,y)\neq(0,0)
}
\]

Sea:

\[
g(u,v)=(x,y)
\]

la inversa local.

Entonces:

\[
J_g(u,v)
=
J_f(x,y)^{-1}.
\]

Así:

\[
\boxed{
J_g(u,v)
=
\frac1{2(x^2+y^2)}
\begin{pmatrix}
x&y\\
-y&x
\end{pmatrix}
}
\]

Además:

\[
f(-x,-y)=f(x,y).
\]

Por tanto:

\[
\boxed{
f\text{ no es inyectiva globalmente}
}
\]

\section{Ejercicio 21}
\subsection{Solución}

Sea:

\[
f(x,y)
=
\left(
x\cos y,
\sen(x-y)
\right).
\]

En:

\[
\left(
\frac{\pi}{2},
\frac{\pi}{2}
\right),
\]

se tiene:

\[
f\left(
\frac{\pi}{2},
\frac{\pi}{2}
\right)
=
(0,0).
\]

La matriz jacobiana es:

\[
J_f(x,y)
=
\begin{pmatrix}
\cos y&-x\sen y\\
\cos(x-y)&-\cos(x-y)
\end{pmatrix}.
\]

En el punto:

\[
J_f\left(
\frac{\pi}{2},
\frac{\pi}{2}
\right)
=
\begin{pmatrix}
0&-\frac{\pi}{2}\\
1&-1
\end{pmatrix}.
\]

Su determinante es:

\[
\frac{\pi}{2}\neq0.
\]

Por tanto, existe una inversa local:

\[
g(u,v).
\]

Además:

\[
J_g(0,0)
=
J_f\left(
\frac{\pi}{2},
\frac{\pi}{2}
\right)^{-1}.
\]

Así:

\[
\boxed{
J_g(0,0)
=
\begin{pmatrix}
-\frac2\pi&1\\[1mm]
-\frac2\pi&0
\end{pmatrix}
}
\]

El jacobiano de \(g\) en ese punto es:

\[
\boxed{
|J_g(0,0)|
=
\frac2\pi
}
\]

\section{Ejercicio 22}
\subsection{Solución}

Sea:

\[
f(x,y,z)
=
\left(
x+xyz,
2y+xy,
z+2x+3z^2
\right).
\]

En:

\[
(0,1,0),
\]

se obtiene:

\[
f(0,1,0)
=
(0,2,0).
\]

La matriz jacobiana es:

\[
J_f(x,y,z)
=
\begin{pmatrix}
1+yz&xz&xy\\
y&2+x&0\\
2&0&1+6z
\end{pmatrix}.
\]

En \((0,1,0)\):

\[
J_f(0,1,0)
=
\begin{pmatrix}
1&0&0\\
1&2&0\\
2&0&1
\end{pmatrix}.
\]

Su determinante es:

\[
|J_f(0,1,0)|=2\neq0.
\]

Por tanto:

\[
\boxed{
f\text{ admite una inversa local }g
\text{ en un entorno de }(0,1,0)
}
\]

Como:

\[
g(0,2,0)=(0,1,0),
\]

se tiene:

\[
J_g(0,2,0)
=
J_f(0,1,0)^{-1}.
\]

Calculando la inversa:

\[
\boxed{
J_g(0,2,0)
=
\begin{pmatrix}
1&0&0\\[1mm]
-\frac12&\frac12&0\\[1mm]
-2&0&1
\end{pmatrix}
}
\]

\end{document}
